\chapter*{Abstract}
% \pdfbookmark[1]{Abstract}{Abstract}
\thispagestyle{empty}
Groundwater (GW) and Surface-water (SW) systems interact seamlessly in reality. They are the two major storage terms in the hydrological cycle that are directly relevant to our civilizations – one hidden beneath the ground and the other very much on display.
\\ \\
In the numerical modelling of our hydrosystems, however, they are often represented by separate numerical models. The reasons for this are multi-fold, as discussed in \ref{intro}. \nameref{intro}.
\\ \\
Although, one must note that the importance of using an integrated numerical model to simulate GW and SW processes becomes crucial in scenarios where the GW-SW interaction fluxes account to a good portion of the total fluxes in the modelling scenario. Often, these scenarios call for the modelling of spatially localized interaction between GW and SW \cite{Winter1998}, like infiltration and exfiltration of GW from a stream for example.
\\ \\
Distributed GW and SW exchange integration in numerical modelling schemes are not as extensively explored as localized GW and SW exchange schemes are \cite{Shen2010}. Increasing interest in intensively studying hydrological fluxes at a regional scale, various interests (including legal, such as the EU Water Framework Directive, 2000) advocating the consideration of GW and SW as a single resource, intensive studies in modelling the fate and transport of pollutants in either GW or SW systems, and special applications are increasing the interest \cite{Fleckenstein2010} in numerical schemes that aim to describe GW and SW in a distributed manner.
\\ \\
In the parent study, it is of key interest to simulate hydrogeological processes of the unsaturated zone, in order to estimate the soil moisture, in a bog wetland study area in order to assess the sustenance of the vegetation it currently hosts. It is suspected that the GW-SW exchange fluxes play a considerable role.
\\ \\
The primary aim of this study is to conceptualize a computationally cheap conceptual model to both, facilitate the coupling of existing GW and SW models in order to to describe a distributed GW-SW exchange, and to simulate the fluxes of water in the unsaturated zone. Additionally, it is also in the interest of this study and the model conceptualized that the model is applied to the study area and that its results are compared and assessed against the measured observations.

\endinput